\lecture{12}{2. Oktober 2025}{Axially Loaded Members}

\section{Axial Load}

\subsection{Saint-Venant's Principle}
In short, Saint-Venant's principle predicts that for an applied force $P$ the distortion of a body will be large locally around the applied force but rather small everywhere else.

\subsection{Elastic Dispalcement of an Axially Loaded Member}
We wish to determine the relative displacement $\delta$ of one end of an axially loaded bar with respect to the other caused by axial loading. We neglect the localized concentrated deformations at the loadings with reference to Saint-Venant's principle. 

\begin{figure} [ht]
  \centering
  \includegraphics[width=0.5\linewidth]{./figures/f12_1.png}
  \caption{Axially loaded bar and differential element.}
  \label{fig:f12_1}
\end{figure}

We consider the bar shown on \Cref{fig:f12_1}. Using the method of sections, we extract a differential element of length $\mathrm{d}x$ and cross sectional area $A(x)$ at position $x$ where the modulus of elasticity is $E(x)$ and the internal axial force is $N(x)$. The stress and strain in the element will then be
\[ 
\sigma = \frac{N(x)}{A(x)} \qquad  \epsilon = \frac{\mathrm{d}\delta}{\mathrm{d}x}
.\]
From Hooke's law the elongation can then be determined as
\[ 
\delta = \int_{0}^{L} \frac{N(x)}{A(x)E(x)}\, \mathrm{d}x
.\]

If the bar is assumed to have constant cross sectional area, and the material is assumed homogeneous and the external force is applied at each end the elongation is simply
\[ 
\delta = \frac{NL}{AE}
.\]
